\documentclass[10.5pt,a4paper]{article}
\usepackage[margin=2cm]{geometry}
\usepackage[T1]{fontenc}
\usepackage{lmodern}
\usepackage{microtype}
\usepackage{hyperref}
\setlength{\parindent}{0pt}\setlength{\parskip}{0.55em}
\pagestyle{empty}
\begin{document}

Quang-Vinh Dang\\
British University Vietnam, Hung Yen, Vietnam\\
vinh.dq4@buv.edu.vn \quad ORCID 0000-0002-3877-8024

\today

The Editors\\
\emph{Computer Vision and Image Understanding}

\textbf{Submission: ``Certified Correction of Object Hallucinations in Vision--Language Models: Guarantees, Gains and Failure Modes under a Pre-specified Protocol''}

Dear Editors,

I submit the above manuscript as an original research article for \emph{Computer Vision and Image Understanding}.

\textbf{Problem.} Vision--language models (VLMs) assert objects that are not in the image, and a wrong answer is costly where they are deployed. Selective prediction lets a model abstain,
and conformal and risk-controlling methods certify the error rate of the answers it does emit. That certificate has a known price: when the model's error rate exceeds the target, any predictor that only emits
or abstains the model's own answer must abstain on a fixed fraction of inputs, whatever its score. The manuscript asks what changes when the predictor may emit a different answer, here the answer of an object detector.

\textbf{Contribution.} I propose Certified-Correction Risk Control (CCRC) for object-existence questions. It accepts the VLM's answer where a correctness score is high, substitutes the answer of an
open-vocabulary detector where the detector's evidence is decisive, abstains otherwise, and certifies the error rate of the emitted set as a whole with a fixed-sequence Clopper--Pearson test. The validity theorem
applies fixed-sequence testing and I claim no new statistical technique. The contribution is the formulation, a protocol in which the proved procedure is the evaluated procedure (folds are unions of whole images, thresholds and repair gate are fixed
before calibration, the start of the test sequence is derived analytically, and the repair gate is fixed on external data, chosen on the fit fold or certified by a union bound), and a measured account of what the mechanism contributes.
The manuscript includes real-image examples of accepted, repaired, abstained and wrongly accepted items.

\textbf{Findings.} On six benchmark--backbone cells (POPE, AMBER; LLaVA-1.5, Qwen2-VL, visual contrastive decoding; OWLv2 as the evidence source) the gain in certified coverage over filtering is $+0.8$ to $+7.3$ points with a gate fixed on external data,
positive in all twelve (cell, risk target) combinations; an outer image-level bootstrap excludes zero in seven. Most of the coverage comes from the correctness score; much of the repaired mass re-emits the model's own answer with detector confirmation;
false-positive errors are corrected only with a large gate; fixed-sequence testing is fragile with small calibration folds; and a small matched-prefix experiment suggests that mid-caption repair increases downstream hallucination.
Negative and inconclusive results are reported with the same prominence as positive ones.

\textbf{Fit with the journal.} The work concerns reliable object perception and language grounding in vision--language models: hallucination, detector evidence, benchmark protocol (including the fact that POPE contributes six questions per image and
splits must be made by image) and risk-controlled deployment. I believe it suits the journal's readership in vision and vision--language understanding.

\textbf{Originality and ethics.} The manuscript has not been published and is not under consideration elsewhere. Earlier versions were submitted to two other journals: one was not accepted after review, and one was declined at the editorial stage as outside that journal's scope. The present
version is substantially revised, with new experiments (extended POPE and AMBER data, an image-grouped protocol, additional analyses) and a restructured text. I have no competing interests to declare and received no specific funding. Generative AI tools were used to support writing and coding, as declared in the manuscript;
all ideas, algorithms and research directions are my own, and I reviewed and edited all content. Per-item model and detector signals, image identifiers and all scripts that generate the tables and figures are publicly available at
\url{https://github.com/vinhqdang/cmvkg} (directory \texttt{conformal\_followup/revision}).

Thank you for considering the manuscript.

Sincerely,

Quang-Vinh Dang

\end{document}
